
% =============================================================================
% Project : THE_PHILOSOPHY_OF_VIABILITY
% Book    : Book I — Ontology of Viability
% Book ID : VS.B01
% Part    : Part I — Identity and Manifestation
% Part ID : VS.B01.P01
% Chapter : Chapter 08 — Comparative Philosophical Examination
% Chapter ID : VS.B01.P01.C08
%
% File :
%   2_MAINMATTER/PART_I_IDENTITY_AND_MANIFESTATION/
%   CHAPTER_08_COMPARATIVE_PHILOSOPHICAL_EXAMINATION/
%   sections/section_13_avicenna.tex
%
% Section :
%   Section 13 — Avicenna (Ibn Sina)
%
% Structural Identifier :
%   VS.B01.P01.C08.S13
%
% Purpose :
%   To reconstruct Avicenna's metaphysics of the existent,
%   essence, quiddity, existence, modality, unity, causality,
%   the Necessary Existent, emanation, and individuation;
%   and to examine the explanatory reach and limits of these
%   positions in relation to the ontology of Being and Identity
%   developed in Book I.
%
% Method :
%   Historical and textual reconstruction;
%   distinction between primary claims and later interpretation;
%   examination of disputed readings and doctrinal qualifications;
%   symmetrical dialectical comparison;
%   strongest challenge and considered response on both sides;
%   bilateral assessment of explanatory costs;
%   substitution test and explicit claim-status classification.
%
% Foundational Restrictions :
%   1. Being and Identity remain analytically distinguishable
%      and ontologically inseparable.
%   2. Analytical Distinction, Ontological Inseparability,
%      and Absence of Ontological Priority remain three
%      independent foundational commitments.
%   3. Identity is not equated with essence, quiddity,
%      determination, unity, singularity, individuation,
%      numerical identity, invariance, structure, composition,
%      material constitution, relational determination,
%      manifestation, persistence, or continuity.
%   4. Avicenna's distinction between essence and existence
%      is not translated into the distinction between
%      Being and Identity.
%   5. The Necessary Existent is not identified with Being
%      as such, Identity, or First Determination.
%   6. The unity and uniqueness of the Necessary Existent
%      are not generalized into a universal thesis that
%      every existent has exactly one identity.
%   7. Causal dependence and emanation are not equated with
%      Expression or Manifestation.
%   8. The hierarchy of intelligences and celestial bodies
%      is not mapped onto the proposed realms and hubs of Book I.
%   9. The Active Intellect is not identified with Light.
%  10. Avicenna's accounts of intellection, perception, and
%      the soul are not made constitutive foundations of
%      Book I's ontology.
%  11. The positive definition and scope of Viability remain
%      open and are not inferred from Avicenna's account
%      of necessity, perfection, or causal dependence.
%
% Claim Classification :
%   Functional or logical form and epistemic or justificatory
%   status are classified independently.
%
% Bibliographical Basis :
%   AvicennaMetaphysics
%   InatiAvicennaRemarksAdmonitions2014
%   LizziniIbnSinaMetaphysicsSEP2026
%   PournamdarAvicennaIndividuation2025
%   BilginTiryakiAvicennaIndividualization2025
%
% Asset Inventory :
%   TL01 — Principal Avicennian claims and comparative limits.
%   DG01 — Modal status and causal dependence.
%   FG   — Not produced; no independent figure is presently
%          justified beyond the table and diagram.
%   IL   — Deferred pending the cumulative visual audit
%          of Chapter 08.
%
% Technical Note :
%   Paragraph destinations use standard hyperref labels,
%   consistent with the existing Chapter 08 sections.
%   The project-wide \InternalRef macro is not defined in
%   the inspected macros.tex; no such macro is assumed here.
%   No new package or global macro is introduced.
%
% =============================================================================


\section{Avicenna}
\label{sec:ch08_s13}

\index[name]{Avicenna}
\index[name]{Ibn Sina}
\index[concept]{Metaphysics}
\index[concept]{Being}
\index[concept]{Existence}
\index[concept]{Existent}
\index[concept]{Essence}
\index[concept]{Quiddity}
\index[concept]{Necessary Existent}
\index[concept]{Possible Existent}
\index[concept]{Necessity}
\index[concept]{Possibility}
\index[concept]{Unity}
\index[concept]{Causality}
\index[concept]{First Cause}
\index[concept]{Emanation}
\index[concept]{Individuation}
\index[concept]{Identity}
\index[concept]{Determination}
\index[concept]{Manifestation}


% =============================================================================
% 1. SCOPE AND COMPARATIVE ORIENTATION
% =============================================================================

\subsection{Scope and Comparative Orientation}
\label{subsec:ch08-avicenna-scope}

% Structural ID: VS.B01.P01.C08.S13.SS01.P001
\phantomsection\label{para:VS.B01.P01.C08.S13.SS01.P001}

Avicenna (Ibn Sina) is a decisive interlocutor for the
ontology of Being and Identity because his metaphysics
connects the analysis of the existent with the questions
of what a thing is, whether it exists necessarily or
possibly, how its existence depends upon a cause, and
how unity and plurality are related to the constitution
of things. His philosophical system does not treat
existence as an isolated fact. It investigates the
modal and causal conditions under which an existent
is what it is and has the existence that it has.

% Structural ID: VS.B01.P01.C08.S13.SS01.P002
\phantomsection\label{para:VS.B01.P01.C08.S13.SS01.P002}

The present examination concentrates on the metaphysical
architecture developed principally in the
\textit{Ilahiyyat} of the \textit{Kitab al-Shifa}
(\textit{The Metaphysics of The Healing}), with
supplementary attention to the metaphysical discussions
in the \textit{al-Najat} and the
\textit{al-Isharat wa-l-Tanbihat}. The central subjects
are the existent as a primary notion, the distinction
between essence and existence, quiddity considered
in itself, the Necessary Existent, possibility and
necessity through another, unity, causal dependence,
emanation, and individuation.

% Structural ID: VS.B01.P01.C08.S13.SS01.P003
\phantomsection\label{para:VS.B01.P01.C08.S13.SS01.P003}

The comparison with Book I must begin by identifying
a fundamental difference in the questions being asked.
Avicenna investigates the relation between what a thing
is and the fact that it exists. Book I investigates
the relation between Being and Identity, maintaining
their analytical distinction, ontological inseparability,
and absence of ontological priority. These are not
alternative expressions of one distinction. The
Avicennian distinction cannot be translated directly
into the foundational distinction adopted in this book.

% Structural ID: VS.B01.P01.C08.S13.SS01.P004
\phantomsection\label{para:VS.B01.P01.C08.S13.SS01.P004}

The principal thesis of this section is that Avicenna
provides an exceptionally developed account of the
modal and causal dependence of existents, together
with a subtle analysis of quiddity and its relation
to existence. His account places a substantial
explanatory demand upon Book I: if quiddity concerns
the reality by virtue of which a thing is what it is,
Book I must explain what additional ontological work
is performed by Identity and why Identity cannot
simply be treated as another name for quiddity.
Nevertheless, the distinction between essence and
existence does not by itself settle the relation
between Being and Identity.

% Structural ID: VS.B01.P01.C08.S13.SS01.P005
\phantomsection\label{para:VS.B01.P01.C08.S13.SS01.P005}

This section does not seek to assimilate Avicenna's
metaphysics to the terminology of the Philosophy
of Viability. Quiddity is not renamed Identity;
existence is not renamed Being; the Necessary
Existent is not identified with First Determination;
and emanation is not treated as another term for
Manifestation. The comparison concerns the
explanatory functions of these concepts and the
questions that remain when their relations are
examined from the standpoint of Book I.


% =============================================================================
% 2. HISTORICAL AND TEXTUAL LOCATION
% =============================================================================

\subsection{Historical and Textual Location}
\label{subsec:ch08-avicenna-historical}

% Structural ID: VS.B01.P01.C08.S13.SS02.P001
\phantomsection\label{para:VS.B01.P01.C08.S13.SS02.P001}

Avicenna was an eleventh-century philosopher whose
work brought Aristotelian philosophy, late-antique
commentary, Arabic Neoplatonic materials, and
Islamic philosophical and theological debates
into a distinctive systematic configuration.
His metaphysics is neither a simple reproduction
of Aristotle nor a mere continuation of Al-Farabi.
It reorders inherited questions and develops
new relations among ontology, theology, cosmology,
and the explanation of causal dependence.

% Structural ID: VS.B01.P01.C08.S13.SS02.P002
\phantomsection\label{para:VS.B01.P01.C08.S13.SS02.P002}

The principal text for this section is the
metaphysical part of the \textit{Kitab al-Shifa},
translated by Michael E. Marmura as
\textit{The Metaphysics of The Healing}. Its
structure is especially important. Book I
introduces metaphysics, the existent, and
the modal division between necessity and
possibility. Book V examines quiddity and
its relation to universality, particularity,
unity, and plurality. Books VI and VIII
develop causal dependence and the Necessary
Existent, while Book IX treats emanation
and the cosmological order. These discussions
must be read in their argumentative contexts
rather than compressed into a single formula.
\citep[Ilahiyyat, Books I, V, VI, VIII--IX]{AvicennaMetaphysics}

% Structural ID: VS.B01.P01.C08.S13.SS02.P003
\phantomsection\label{para:VS.B01.P01.C08.S13.SS02.P003}

The \textit{al-Isharat wa-l-Tanbihat} is also
relevant because it presents mature formulations
of existence, its causes, creation, and the
arrangement of existents. Shams C. Inati's
annotated translation of its Physics and
Metaphysics provides a further textual basis
for examining Avicenna's treatment of causal
dependence and the relation between the
Necessary Existent and what depends upon it.
The use of this work is supplementary: it
does not replace the more extensive analysis
of the \textit{Ilahiyyat} in the
\textit{Kitab al-Shifa}.
\citep[Metaphysics, Fourth--Sixth Classes]{InatiAvicennaRemarksAdmonitions2014}

% Structural ID: VS.B01.P01.C08.S13.SS02.P004
\phantomsection\label{para:VS.B01.P01.C08.S13.SS02.P004}

Avicenna's different philosophical works must not be treated
as a single undifferentiated textual witness. The present
reconstruction uses the \textit{Kitab al-Shifa} as its principal
primary source and the metaphysical discussions in the
\textit{al-Isharat wa-l-Tanbihat} as supplementary evidence.
Where a claim depends specifically on another work, its
textual basis must be established independently rather than
inferred from the similarity of its subject matter or from
a later summary. This qualification is necessary to prevent
the construction of an artificially uniform doctrine from
texts whose arguments, purposes, and expositions may differ.
\citep[Metaphysics, Books I--II]{AvicennaMetaphysics}

% Structural ID: VS.B01.P01.C08.S13.SS02.P005
\phantomsection\label{para:VS.B01.P01.C08.S13.SS02.P005}

The historical relation between Al-Farabi
and Avicenna is especially important for
the preceding section. Both investigate
Being, quiddity, unity, the First Principle,
and causal dependence. Avicenna's treatment,
however, develops a more systematic modal
analysis of the relation between a thing's
quiddity and its existence. The later
formulation must not be projected backward
onto Al-Farabi, just as Al-Farabi's
terminology must not be assumed to contain
all the distinctions developed by Avicenna.
The comparison therefore concerns both
continuity of problems and transformation
of explanatory structure.
\footnote{For the historical relation between
Avicenna's metaphysics, Al-Farabi, Aristotelian
commentary, and Arabic Neoplatonic materials,
see Olga Lizzini, ``Ibn Sina's Metaphysics,''
\textit{The Stanford Encyclopedia of Philosophy},
substantive revision 9 June 2026. The account
is used here as secondary historical and
interpretive guidance, not as a substitute
for the primary texts. \cite{LizziniIbnSinaMetaphysicsSEP2026}}


% =============================================================================
% 3. METAPHYSICS AS A SCIENCE
% =============================================================================

\subsection{Metaphysics as a Science}
\label{subsec:ch08-avicenna-metaphysics-science}

% Structural ID: VS.B01.P01.C08.S13.SS03.P001
\phantomsection\label{para:VS.B01.P01.C08.S13.SS03.P001}

Avicenna defines metaphysics as the science
of the existent insofar as it is existent.
Its subject matter is not a particular
class of entities, such as bodies, numbers,
or living things, but the existent considered
without restriction to one particular
domain or category. The universality of
metaphysics therefore concerns the scope
of its subject matter and the principles
it investigates. It does not mean that
all existents are identical or that their
differences can be eliminated by a single
abstract definition.
\citep[Ilahiyyat, I.1--2]{AvicennaMetaphysics}

% Structural ID: VS.B01.P01.C08.S13.SS03.P002
\phantomsection\label{para:VS.B01.P01.C08.S13.SS03.P002}

Avicenna distinguishes the subject matter
of a science from the questions whose
answers the science must establish.
Metaphysics begins from the existent as
its subject, while the existence of the
First Principle is among the matters that
it must demonstrate. This distinction
allows metaphysics to be both ontology
and theology without reducing its entire
subject matter to the investigation of
a divine entity. The First Principle is
investigated within metaphysics because
it is existent, while its status as the
uncaused principle of caused existents
constitutes a distinct question.
\citep[Ilahiyyat, I.1--3]{AvicennaMetaphysics}

% Structural ID: VS.B01.P01.C08.S13.SS03.P003
\phantomsection\label{para:VS.B01.P01.C08.S13.SS03.P003}

The distinction is philosophically important
because it prevents two reductions. First,
metaphysics is not merely a theological
science whose subject is God and whose
remaining questions are secondary
consequences. Second, it is not merely
a neutral classification of entities
that can avoid the question of their
ultimate causal explanation. Avicenna
connects the universal investigation
of the existent with the demonstration
of a First Principle, while preserving
the difference between the subject
matter and the object of demonstration.
\citep[Ilahiyyat, I.1--3]{AvicennaMetaphysics}

% Structural ID: VS.B01.P01.C08.S13.SS03.P004
\phantomsection\label{para:VS.B01.P01.C08.S13.SS03.P004}

For Book I, this structure raises a
methodological question. An ontology
of Being and Identity must clarify
whether its foundational distinctions
are sufficient to articulate the
existent as such, or whether they
presuppose further principles concerning
the cause and modal status of existence.
Avicenna's approach shows that an
ontology may require both a universal
account of the existent and an argument
concerning the dependence of particular
existents. It does not establish that
the specific causal architecture of
Avicenna must be adopted by Book I.

% Structural ID: VS.B01.P01.C08.S13.SS03.P005
\phantomsection\label{para:VS.B01.P01.C08.S13.SS03.P005}

The fact that Avicenna describes primary
concepts as being impressed in the soul
does not make the existence of what they
concern dependent upon cognition. The
epistemic or conceptual priority of a
notion within inquiry must be distinguished
from the ontological status of its
referent. Book I may therefore examine
Avicenna's account of primary concepts
without adopting a cognitive subject
as the constitutive ground of Being
or Identity.
\footnote{The distinction between the
primary conceptual character of the
existent and the ontological independence
of what exists is discussed in Lizzini,
``Ibn Sina's Metaphysics,'' sections 2--3.
\cite{LizziniIbnSinaMetaphysicsSEP2026}}


% =============================================================================
% 4. THE EXISTENT AS A PRIMARY NOTION
% =============================================================================

\subsection{The Existent as a Primary Notion}
\label{subsec:ch08-avicenna-existent-primary}

% Structural ID: VS.B01.P01.C08.S13.SS04.P001
\phantomsection\label{para:VS.B01.P01.C08.S13.SS04.P001}

Avicenna treats the existent (\textit{al-mawjud})
as a primary notion that cannot be defined
through a more basic concept. Any attempt
to define it would already employ the
notion of what is or what is not. The
existent is therefore not introduced
as one genus among others, nor as a
particular entity that could be classified
under a more universal category. It is
a starting point of metaphysical inquiry.
\citep[Ilahiyyat, I.2 and I.5]{AvicennaMetaphysics}

% Structural ID: VS.B01.P01.C08.S13.SS04.P002
\phantomsection\label{para:VS.B01.P01.C08.S13.SS04.P002}

The primary character of the existent
does not imply that every use of the
term has exactly the same function.
Avicenna distinguishes the existent
from related notions such as thing
(\textit{shay'}), whatness, and the
reality by virtue of which a thing
is what it is. The terms are closely
related and may be extensionally
coincident in particular contexts,
but their conceptual functions are
not thereby identical. The analysis
must distinguish what is affirmed
when something is said to exist from
what is specified when its thingness
or quiddity is considered.
\citep[Ilahiyyat, I.5]{AvicennaMetaphysics}

% Structural ID: VS.B01.P01.C08.S13.SS04.P003
\phantomsection\label{para:VS.B01.P01.C08.S13.SS04.P003}

This distinction prevents a common
misreading. The fact that the existent
is a primary notion does not mean
that existence alone supplies a
complete account of the thing that
exists. Conversely, the fact that
a thing has a quiddity does not
mean that its existence follows
from the definition of that quiddity.
The question of what a thing is
and the question of whether it exists
are related but analytically
distinguishable.
\citep[Ilahiyyat, I.5--6]{AvicennaMetaphysics}

% Structural ID: VS.B01.P01.C08.S13.SS04.P004
\phantomsection\label{para:VS.B01.P01.C08.S13.SS04.P004}

The comparison with Book I must
preserve this distinction without
translating it into the foundational
relation between Being and Identity.
Avicenna's distinction concerns
the existent and its quiddity,
or the affirmation of existence
and the reality specified by
whatness. Book I's distinction
concerns Being and Identity.
The terms overlap in the questions
they make possible, but their
inferential roles differ. No
equivalence follows merely from
the fact that both analyses
distinguish two aspects of an
existent.


% =============================================================================
% 5. ESSENCE, QUIDDITY, AND EXISTENCE
% =============================================================================

\subsection{Essence, Quiddity, and Existence}
\label{subsec:ch08-avicenna-essence-existence}

% Structural ID: VS.B01.P01.C08.S13.SS05.P001
\phantomsection\label{para:VS.B01.P01.C08.S13.SS05.P001}

Avicenna's distinction between essence
and existence is not adequately described
as the difference between a thing and
an additional property subsequently
attached to it. The distinction concerns
the relation between what a thing is
and the actuality through which it
exists. A quiddity may be understood
without thereby establishing that
the thing possessing that quiddity
exists in external reality. The
analysis of whatness therefore does
not, by itself, demonstrate the
existence of its subject.
\citep[Ilahiyyat, I.5--6; V.1]{AvicennaMetaphysics}

% Structural ID: VS.B01.P01.C08.S13.SS05.P002
\phantomsection\label{para:VS.B01.P01.C08.S13.SS05.P002}

The distinction is not a claim that
quiddity is an independently existing
object waiting to receive existence.
Avicenna's account does not require
a second entity called essence to
be placed alongside an existent
entity called existence. Such a
picture would turn the distinction
into an unnecessary multiplication
of entities and would generate
the further question of how the
essence itself exists. The
distinction is metaphysical and
explanatory, but it must not be
reified into two independently
existing things.
\citep[Ilahiyyat, V.1 and VIII.4]{AvicennaMetaphysics}

% Structural ID: VS.B01.P01.C08.S13.SS05.P003
\phantomsection\label{para:VS.B01.P01.C08.S13.SS05.P003}

In contingent existents, the quiddity
does not entail its own existence.
Its existence is dependent upon
a cause. In the Necessary Existent,
by contrast, there is no quiddity
distinct from existence in the
same manner as in caused existents.
The distinction therefore has
a modal and causal function: it
helps explain why the existence
of a possible thing requires
a cause, while the Necessary
Existent does not receive its
existence from another.
\citep[Ilahiyyat, I.6--7; VIII.4]{AvicennaMetaphysics}

% Structural ID: VS.B01.P01.C08.S13.SS05.P004
\phantomsection\label{para:VS.B01.P01.C08.S13.SS05.P004}

The relation between essence and
existence is thus more complex
than an abstract contrast between
a definition and an act of being.
Avicenna's analysis concerns the
modal status of a quiddity and
its relation to the causal conditions
under which it exists. The
distinction is not merely epistemic:
it is intended to explain an
ontological dependence that
cannot be inferred from the
quiddity considered by itself.
\citep[Ilahiyyat, I.5--7; VIII.1--4]{AvicennaMetaphysics}

% Structural ID: VS.B01.P01.C08.S13.SS05.P005
\phantomsection\label{para:VS.B01.P01.C08.S13.SS05.P005}

The comparison with Book I must
not begin by assuming that the
Avicennian distinction is either
identical with or opposed to the
distinction between Being and
Identity. Book I does not claim
that Being and Identity are two
independently existing components.
It maintains that they are
analytically distinguishable
while ontologically inseparable.
This position does not, by itself,
deny that an existent may be
analysed through quiddity and
existence. The precise relation
between the two pairs of concepts
remains to be examined.

% Structural ID: VS.B01.P01.C08.S13.SS05.P006
\phantomsection\label{para:VS.B01.P01.C08.S13.SS05.P006}

The decisive question is therefore
not whether Book I should reject
the distinction between essence
and existence. It is whether
that distinction exhausts the
ontological problem of Identity.
If the quiddity specifies what
a thing is, the question remains
whether it also accounts for
the complete ontological status
of the existent as the existent
that it is, including the
conditions of its individuality,
its relations, and its possible
continuity through transformation.


% =============================================================================
% 6. QUIDDITY CONSIDERED IN ITSELF
% =============================================================================

\subsection{Quiddity Considered in Itself}
\label{subsec:ch08-avicenna-quiddity-itself}

% Structural ID: VS.B01.P01.C08.S13.SS06.P001
\phantomsection\label{para:VS.B01.P01.C08.S13.SS06.P001}

One of the most influential features
of Avicenna's metaphysics is the
analysis of quiddity considered
in itself. A quiddity considered
precisely as the quiddity that
it is does not, in that consideration,
include the conditions of its
existence as either universal
or particular. The example of
horseness is instructive: horseness
as such is neither universal
nor particular. Universality and
particularity arise when the
quiddity is considered in
different modes of existence
and under their accompanying
conditions.
\citep[Ilahiyyat, V.1--2]{AvicennaMetaphysics}

% Structural ID: VS.B01.P01.C08.S13.SS06.P002
\phantomsection\label{para:VS.B01.P01.C08.S13.SS06.P002}

The doctrine must not be interpreted
as the claim that a quiddity exists
as a separate object in a third
realm before it becomes instantiated.
Avicenna's account does not require
an independently subsisting universal
or a separately existing essence.
Rather, it distinguishes the
quiddity considered in itself
from the conditions accompanying
it when it exists in the mind
or in external individuals.
The exact ontological interpretation
of this distinction remains
a subject of scholarly debate.
\citep[Ilahiyyat, V.1--2]{AvicennaMetaphysics}

% Structural ID: VS.B01.P01.C08.S13.SS06.P003
\phantomsection\label{para:VS.B01.P01.C08.S13.SS06.P003}

The importance of the analysis
for the problem of Identity is
substantial. If quiddity is not
itself universal or particular,
then the relation between a
common whatness and its individual
instances requires further
explanation. The quiddity alone
does not specify whether it is
present as a universal in thought
or as a particular in external
reality. The conditions accompanying
its existence therefore become
relevant to the analysis of
individuality and plurality.
\citep[Ilahiyyat, V.1--2]{AvicennaMetaphysics}

% Structural ID: VS.B01.P01.C08.S13.SS06.P004
\phantomsection\label{para:VS.B01.P01.C08.S13.SS06.P004}

This account creates a serious
comparative question for Book I.
If quiddity is the reality by
virtue of which a thing is what
it is, and if it is considered
without presupposing universality
or particularity, then Identity
cannot be defended merely by
repeating that an existent is
the existent that it is. Book I
must explain whether Identity
is related to quiddity as a
distinct ontological dimension,
whether quiddity is one aspect
of the constitution of Identity,
or whether the relation requires
a more complex account that
cannot yet be formulated.
\citep[Ilahiyyat, I.5 and V.1]{AvicennaMetaphysics}

% Structural ID: VS.B01.P01.C08.S13.SS06.P005
\phantomsection\label{para:VS.B01.P01.C08.S13.SS06.P005}

The present section does not
resolve this question by declaring
quiddity to be merely conceptual.
Nor does it identify quiddity
with Identity. Both moves would
avoid the central problem rather
than answer it. The issue is
whether the account of whatness
provides a complete explanation
of the existent's identity or
whether further ontological
distinctions remain necessary.


% =============================================================================
% 7. THE NECESSARY EXISTENT IN ITSELF
% =============================================================================

\subsection{The Necessary Existent in Itself}
\label{subsec:ch08-avicenna-necessary-existent}

% Structural ID: VS.B01.P01.C08.S13.SS07.P001
\phantomsection\label{para:VS.B01.P01.C08.S13.SS07.P001}

Avicenna distinguishes what is
necessary in itself from what
is possible in itself. The
Necessary Existent in itself
does not receive existence from
a cause. Its existence is not
dependent upon a condition
outside itself, and its status
cannot be explained through
a prior cause. The Necessary
Existent is therefore uncaused
and is the ultimate principle
within Avicenna's account of
causal dependence.
\citep[Ilahiyyat, I.6--7]{AvicennaMetaphysics}

% Structural ID: VS.B01.P01.C08.S13.SS07.P002
\phantomsection\label{para:VS.B01.P01.C08.S13.SS07.P002}

The Necessary Existent is also
characterized through unity,
simplicity, and uniqueness.
Avicenna argues that there
cannot be two independently
necessary existents of the
same absolute status. If two
such existents were distinguished
by an essential difference,
neither would be necessary
in itself in the required
sense; if they were distinguished
by an external or accidental
condition, their distinction
would introduce dependence
upon something other than
themselves. The argument
supports the uniqueness of
the Necessary Existent within
Avicenna's system.
\citep[Ilahiyyat, I.7 and VIII.4]{AvicennaMetaphysics}

% Structural ID: VS.B01.P01.C08.S13.SS07.P003
\phantomsection\label{para:VS.B01.P01.C08.S13.SS07.P003}

The unity of the Necessary
Existent must not be confused
with a universal account of
the unity of every existent.
Avicenna's argument concerns
the unique status of that
which is necessary in itself
and independent of every
cause. It does not establish
that every caused existent
is simple, that every
aggregate has one identity,
or that every existent must
possess exactly one identity.
The uniqueness of the First
Principle is a conclusion
within a particular modal
and causal argument.
\citep[Ilahiyyat, I.6--7]{AvicennaMetaphysics}

% Structural ID: VS.B01.P01.C08.S13.SS07.P004
\phantomsection\label{para:VS.B01.P01.C08.S13.SS07.P004}

The Necessary Existent is not
identified with Being as such.
Avicenna investigates the
Necessary Existent as the
uncaused principle whose
existence is necessary in
itself. The existent as such,
however, is the universal
subject of metaphysics and
is not reducible to the
First Principle. The First
Principle is therefore not
equivalent to the totality
of Being, to Identity, or
to First Determination in
the proposed architecture
of Book I.
\citep[Ilahiyyat, I.1--3 and VIII.4]{AvicennaMetaphysics}


% =============================================================================
% 8. POSSIBILITY AND NECESSITY THROUGH ANOTHER
% =============================================================================

\subsection{Possibility and Necessity through Another}
\label{subsec:ch08-avicenna-possibility-necessity}

% Structural ID: VS.B01.P01.C08.S13.SS08.P001
\phantomsection\label{para:VS.B01.P01.C08.S13.SS08.P001}

A possible existent considered
in itself does not contain
within its quiddity a sufficient
reason for its existence.
Its quiddity does not make
its existence necessary,
and it does not make its
non-existence impossible.
The existent is therefore
possible in itself. This
possibility is not merely
an expression of human
uncertainty; it identifies
a modal status that Avicenna
regards as belonging to
the thing considered in
itself.
\citep[Ilahiyyat, I.6]{AvicennaMetaphysics}

% Structural ID: VS.B01.P01.C08.S13.SS08.P002
\phantomsection\label{para:VS.B01.P01.C08.S13.SS08.P002}

When the possible existent
is related to its complete
cause, its existence becomes
necessary through another.
This does not mean that
the existent ceases to be
possible in itself. The
two descriptions concern
different respects: considered
in itself, the thing is
possible; considered in
relation to the cause
that makes it exist, it
is necessary through
another. The modal
distinction therefore
preserves the difference
between what a thing is
in itself and the causal
conditions under which
it exists.
\citep[Ilahiyyat, I.6 and VIII.1--3]{AvicennaMetaphysics}

% Structural ID: VS.B01.P01.C08.S13.SS08.P003
\phantomsection\label{para:VS.B01.P01.C08.S13.SS08.P003}

The distinction between
necessity in itself and
necessity through another
is central to Avicenna's
account of dependence.
A caused existent does
not become uncaused merely
because its existence is
necessary when its complete
cause is present. Its
necessity remains dependent
and relational. The
Necessary Existent, by
contrast, has no such
causal relation as the
ground of its existence.
\citep[Ilahiyyat, I.6--7; VIII.1--4]{AvicennaMetaphysics}

% Structural ID: VS.B01.P01.C08.S13.SS08.P004
\phantomsection\label{para:VS.B01.P01.C08.S13.SS08.P004}

This modal account is important
for Book I because it
distinguishes the ontological
status of an existent from
the conditions under which
its existence is explained.
Nevertheless, necessity
through another must not
be equated with persistence,
stability, invariance, or
the preservation of Identity.
A thing's being causally
necessitated does not,
by itself, establish
whether its Identity
continues through every
transformation or whether
its determinations remain
invariant under a specified
domain of change.

% Structural ID: VS.B01.P01.C08.S13.SS08.P005
\phantomsection\label{para:VS.B01.P01.C08.S13.SS08.P005}

The diagram inserted below
represents the modal and
causal relations in a
restricted explanatory
form. It does not depict
a temporal transition
from possibility to
necessity. The same
caused existent is
possible in itself and
necessary through another
when considered under
the two relevant respects.
The causal relation
explains its existence;
it does not provide
a complete account
of its Identity.


% =============================================================================
% Project : THE_PHILOSOPHY_OF_VIABILITY
% Book    : Book I — Ontology of Viability
% Book ID : VS.B01
% Part    : Part I — Identity and Manifestation
% Part ID : VS.B01.P01
% Chapter : Chapter 08 — Comparative Philosophical Examination
% Chapter ID : VS.B01.P01.C08
%
% File :
%   4_ASSETS/diagrams/
%   VS.B01.P01.C08.S13.SS08.P005.DG01.avicenna_modal_causal_dependence.tex
%
% Asset :
%   DG01 — Avicenna's Modal and Causal Dependence
%
% Structural Identifier :
%   VS.B01.P01.C08.S13.SS08.P005.DG01
%
% Purpose :
%   To distinguish the possible existent considered in itself
%   from the same existent considered as necessary through
%   its complete cause, and to indicate the explanatory
%   dependence of caused existents upon the First Principle.
%
% Dependencies :
%   tikz, positioning, arrows.meta.
%   These are already loaded by the Book I preamble.
%
% Philosophical and Representational Limits :
%   1. The diagram is schematic and not an exhaustive map
%      of Avicenna's cosmology.
%   2. The two modal descriptions concern the same caused
%      existent under different respects; they are not
%      successive temporal stages.
%   3. The arrows indicate causal or explanatory dependence,
%      not temporal succession.
%   4. The causal series is simplified for explanatory clarity.
%   5. The diagram does not imply that the First Principle
%      is identical with Being as such or First Determination.
%   6. The diagram does not map Avicenna's cosmology onto
%      the proposed realms and hubs of Book I.
%   7. Causal dependence is not equated with Manifestation.
%
% =============================================================================

\begin{figure}[htbp]
    \centering

    \begin{tikzpicture}[
        node distance=0.78cm,
        principle/.style={
            draw=ViabilityGold,
            fill=ViabilityPaleGold,
            rounded corners=2pt,
            line width=0.8pt,
            text width=0.70\linewidth,
            align=center,
            inner sep=7pt
        },
        mediation/.style={
            draw=ViabilityBlue,
            fill=ViabilityPaleBlue,
            rounded corners=2pt,
            line width=0.7pt,
            text width=0.70\linewidth,
            align=center,
            inner sep=7pt
        },
        caused/.style={
            draw=ViabilityGreen,
            fill=ViabilityPaleGreen,
            rounded corners=2pt,
            line width=0.8pt,
            text width=0.82\linewidth,
            align=center,
            inner sep=8pt
        },
        dependence/.style={
            -{Stealth[length=2.2mm]},
            draw=ViabilityDarkBlue,
            line width=0.85pt
        },
        relationlabel/.style={
            fill=white,
            text=ViabilityDarkBlue,
            font=\footnotesize,
            inner sep=2pt,
            align=center
        }
    ]

    \node[principle] (first) {
        \textbf{Necessary Existent in Itself}\\[3pt]
        Uncaused; independent; not necessitated by another
    };

    \node[mediation, below=of first] (series) {
        \textbf{Causal Order}\\[3pt]
        Intermediate causes and relations of dependence
    };

    \node[mediation, below=of series] (cause) {
        \textbf{Complete Cause of a Caused Existent}\\[3pt]
        The causal conditions sufficient for its existence
    };

    \node[caused, below=of cause] (effect) {
        \textbf{One Caused Existent Considered under Two Respects}\\[7pt]

        \begin{tabular}{c}
            \textbf{In itself:}\\
            Possible Existent
            \\[5pt]
            \textbf{In relation to its complete cause:}\\
            Necessary through Another
        \end{tabular}
    };

    \draw[dependence]
        (first) --
        node[relationlabel, right] {ultimate causal\\grounding}
        (series);

    \draw[dependence]
        (series) --
        node[relationlabel, right] {mediated causal\\dependence}
        (cause);

    \draw[dependence]
        (cause) --
        node[relationlabel, right] {necessitation\\through another}
        (effect);

    \end{tikzpicture}

    \caption[
        Avicenna's distinction between possibility in itself
        and necessity through another
    ]{
        Avicenna's modal and causal account of a caused existent.
        The possible existent remains possible when considered
        in itself, while its existence is necessary through
        another when considered in relation to its complete
        cause. The causal order is represented schematically;
        the diagram does not depict a temporal sequence.
    }

    \label{fig:VS.B01.P01.C08.S13.SS08.P005.DG01}

\end{figure}

% =============================================================================
% END OF DIAGRAM DG01
% =============================================================================



% =============================================================================
% 9. UNITY, ONENESS, AND THE MANY
% =============================================================================

\subsection{Unity, Oneness, and the Many}
\label{subsec:ch08-avicenna-unity}

% Structural ID: VS.B01.P01.C08.S13.SS09.P001
\phantomsection\label{para:VS.B01.P01.C08.S13.SS09.P001}

Unity is a central metaphysical
notion in Avicenna's account.
The existent and the one
are primary notions, and
the one is closely related
to the existent. Yet
their conceptual functions
are not identical. The
analysis of what it means
for something to be one
must distinguish the
different senses in which
unity is predicated and
the different kinds of
entities to which those
predications apply.
\citep[Ilahiyyat, I.5 and III.2--3]{AvicennaMetaphysics}

% Structural ID: VS.B01.P01.C08.S13.SS09.P002
\phantomsection\label{para:VS.B01.P01.C08.S13.SS09.P002}

Avicenna's analysis does
not permit unity to be
treated as a simple
synonym for Identity.
A thing may be considered
one in a particular
respect without this
consideration supplying
a complete account of
what makes it the thing
that it is. Likewise,
the unity of a quiddity,
the unity of an individual,
and the absolute unity
of the Necessary Existent
do not necessarily perform
the same explanatory
function.
\citep[Ilahiyyat, III.2--3 and V.1]{AvicennaMetaphysics}

% Structural ID: VS.B01.P01.C08.S13.SS09.P003
\phantomsection\label{para:VS.B01.P01.C08.S13.SS09.P003}

The analysis of quiddity
considered in itself
also complicates the
relation between unity
and plurality. Quiddity
as such is not universal
or particular, and it
is not, in that restricted
consideration, one or
many. Unity and plurality
become relevant when
the quiddity is considered
under the conditions
accompanying its existence.
The distinction is therefore
not a denial of the
reality of individuals
or of plurality; it
is an analysis of the
conditions under which
quiddity is considered.
\citep[Ilahiyyat, V.1--2]{AvicennaMetaphysics}

% Structural ID: VS.B01.P01.C08.S13.SS09.P004
\phantomsection\label{para:VS.B01.P01.C08.S13.SS09.P004}

For Book I, the important
point is that unity and
Identity remain distinct
problems. The unity of
a source does not establish
the numerical singularity
of every effect. The
unity of an aggregate
does not establish that
the aggregate possesses
exactly one identity.
And the plurality of
individuals sharing
a quiddity does not
settle whether one
existent may involve
multiple identities.
These questions require
independent arguments
and cannot be answered
by transferring the
unity of the Necessary
Existent to every other
level of reality.

% Structural ID: VS.B01.P01.C08.S13.SS09.P005
\phantomsection\label{para:VS.B01.P01.C08.S13.SS09.P005}

The comparative inquiry
must also distinguish
the unity of a quiddity
from the unity of its
instances. A common
quiddity may be predicated
of multiple individuals,
but the commonality of
whatness does not make
those individuals numerically
identical. Conversely,
the fact that two individuals
are numerically distinct
does not establish that
their quiddities are
different in every respect.
The relation between
quiddity, unity, plurality,
and Identity remains
more complex than a
single principle of
oneness can resolve.


% =============================================================================
% 10. CAUSALITY AND THE FIRST PRINCIPLE
% =============================================================================

\subsection{Causality and the First Principle}
\label{subsec:ch08-avicenna-causality}

% Structural ID: VS.B01.P01.C08.S13.SS10.P001
\phantomsection\label{para:VS.B01.P01.C08.S13.SS10.P001}

Avicenna's account of causality
is not confined to the
production of motion or
to the alteration of
an already existing
entity. In metaphysics,
the efficient cause
is fundamentally a
cause of existence.
The question is not
only what initiates
a process but what
accounts for the
existence of a caused
existent. This distinction
allows Avicenna to
connect causal analysis
with the modal status
of the possible existent.
\citep[Ilahiyyat, VI.1--2 and VIII.1--3]{AvicennaMetaphysics}

% Structural ID: VS.B01.P01.C08.S13.SS10.P002
\phantomsection\label{para:VS.B01.P01.C08.S13.SS10.P002}

A causal series includes
a first cause, intermediate
causes, and effects.
An intermediate cause
may itself be an effect
in relation to a prior
cause while functioning
as a cause in relation
to what follows. The
presence of intermediate
causes does not eliminate
the need for an ultimate
explanation of the
existence of the series.
Avicenna's argument
therefore concerns the
dependence of caused
existents and the
conditions under which
a causal order can
be explained.
\citep[Ilahiyyat, VIII.1--3]{AvicennaMetaphysics}

% Structural ID: VS.B01.P01.C08.S13.SS10.P003
\phantomsection\label{para:VS.B01.P01.C08.S13.SS10.P003}

The First Principle
is not simply the first
member of a series
whose other members
are arranged after
it in temporal succession.
Its priority is causal
and ontological. The
causal series may be
eternal in temporal
duration while remaining
dependent upon a first
principle in the order
of explanation. The
distinction between
temporal beginning
and causal dependence
is therefore indispensable
to the reconstruction
of Avicenna's position.
\citep[Ilahiyyat, VIII.1--3 and IX.1]{AvicennaMetaphysics}

% Structural ID: VS.B01.P01.C08.S13.SS10.P004
\phantomsection\label{para:VS.B01.P01.C08.S13.SS10.P004}

The causal argument
must not be reduced
to the assertion that
an infinite temporal
past is impossible.
Avicenna's concern
is with the explanatory
dependence of effects
and with the impossibility
of accounting for
a caused existent
solely through a series
whose members are
themselves dependent
without an ultimate
principle. The distinction
between an eternal
succession and a
dependent causal
order is essential
to a fair comparison.
\footnote{For the structure
of Avicenna's causal
argument and the debate
over its relation to
the demonstration of
the First Principle,
see Lizzini, ``Ibn Sina's
Metaphysics,'' sections
4--5. The argument
is not presented here
as a simple proof
against every possible
form of infinite series.
\cite{LizziniIbnSinaMetaphysicsSEP2026}}

% Structural ID: VS.B01.P01.C08.S13.SS10.P005
\phantomsection\label{para:VS.B01.P01.C08.S13.SS10.P005}

The relation between
causal dependence and
Identity remains a
distinct problem. An
account of why an
existent exists does
not automatically
provide a complete
account of what
constitutes its Identity.
Nevertheless, the
causal conditions of
existence may be
relevant to the
constitution, individuation,
or preservation of
an existent. Their
precise role must
be established rather
than presumed. Book I
cannot dismiss causal
dependence as external
to every ontological
question about Identity.


% =============================================================================
% 11. EMANATION AND THE COSMOLOGICAL ORDER
% =============================================================================

\subsection{Emanation and the Cosmological Order}
\label{subsec:ch08-avicenna-emanation-cosmology}

% Structural ID: VS.B01.P01.C08.S13.SS11.P001
\phantomsection\label{para:VS.B01.P01.C08.S13.SS11.P001}

Avicenna's account of
emanation (\textit{fayd})
explains how the existence
of dependent beings
proceeds from the
Necessary Existent.
The First Principle
is not altered by
the existence of its
effects, and the
causal relation does
not require a temporal
beginning. Emanation
therefore expresses
an ontological dependence
that is compatible
with the eternity
of the caused order.
\citep[Ilahiyyat, IX.1--4]{AvicennaMetaphysics}

% Structural ID: VS.B01.P01.C08.S13.SS11.P002
\phantomsection\label{para:VS.B01.P01.C08.S13.SS11.P002}

The first caused intelligence
occupies a distinctive
position in this order.
Avicenna's account explains
how multiplicity follows
from the First through
mediated causal relations.
The first intelligence
is not absolutely simple
in the manner of the
First Principle. Its
own modal constitution
and acts of intellection
are involved in the
derivation of further
intelligences and
celestial entities.
The familiar formula
that from the One,
insofar as it is one,
only one proceeds
belongs to this
specific cosmological
argument and must
not be generalized
without qualification.
\citep[Ilahiyyat, IX.4--5]{AvicennaMetaphysics}

% Structural ID: VS.B01.P01.C08.S13.SS11.P003
\phantomsection\label{para:VS.B01.P01.C08.S13.SS11.P003}

The cosmological structure
includes separate intelligences,
celestial souls and
bodies, and the
sublunary order.
The details of the
series and its relation
to the celestial spheres
must be reconstructed
from the relevant
texts. The number
and arrangement of
the intelligences
should not be treated
as an unqualified
constant across every
Avicennian work or
every later interpretation.
\citep[Ilahiyyat, IX.3--5]{AvicennaMetaphysics}

% Structural ID: VS.B01.P01.C08.S13.SS11.P004
\phantomsection\label{para:VS.B01.P01.C08.S13.SS11.P004}

The Active Intellect
has a significant
role in the explanation
of the sublunary
order and in Avicenna's
account of human
intellectual activity.
Its function belongs
to a specific cosmological
and psychological
architecture. It must
not be identified
with the Ontological
Light investigated
in Part III of Book I.
The fact that both
may be discussed
in relation to
intelligibility or
intellectual activity
does not establish
that they perform
the same ontological
function.
\citep[Ilahiyyat, IX.5 and X.1]{AvicennaMetaphysics}

% Structural ID: VS.B01.P01.C08.S13.SS11.P005
\phantomsection\label{para:VS.B01.P01.C08.S13.SS11.P005}

The cosmological order
is also not equivalent
to the proposed
four realms and
three hubs of Book I.
The latter are
an open and developing
ontological architecture,
not a translation
of the celestial
hierarchy. The order
of exposition in
Book I does not
establish a linear
causal hierarchy,
and the arrangement
of Avicenna's intelligences
does not provide
a ready-made map
of the proposed
domains of Being.


% =============================================================================
% 12. INDIVIDUATION AND THE INDIVIDUAL
% =============================================================================

\subsection{Individuation and the Individual}
\label{subsec:ch08-avicenna-individuation}

% Structural ID: VS.B01.P01.C08.S13.SS12.P001
\phantomsection\label{para:VS.B01.P01.C08.S13.SS12.P001}

The problem of individuation
asks how an individual
is this individual
rather than another
instance of the same
species or quiddity.
Avicenna's discussions
do not yield a single
simple principle that
can be applied without
qualification to every
kind of existent.
His works consider
several possible
individuative factors,
including matter,
accidental features,
form, position,
place, and existence.
Their roles differ
according to the
context and the
kind of entity
under examination.
\citep[Ilahiyyat, V.1--2]{AvicennaMetaphysics}

% Structural ID: VS.B01.P01.C08.S13.SS12.P002
\phantomsection\label{para:VS.B01.P01.C08.S13.SS12.P002}

A distinction between
\textit{tashakhkhus}
and \textit{tamayyuz}
is important in recent
interpretive work.
The former concerns
the individual as
individuated or
possessing its
individuality; the
latter concerns
distinction or
being differentiated
from another. These
terms should not
be treated as
interchangeable. The
question of what
makes an individual
the individual it
is is not exhausted
by the question
of how it can
be distinguished
from another.
\citep{PournamdarAvicennaIndividuation2025}

% Structural ID: VS.B01.P01.C08.S13.SS12.P003
\phantomsection\label{para:VS.B01.P01.C08.S13.SS12.P003}

Recent interpretations
also emphasize the
relation between
quiddity considered
in itself, individual
existence, and the
field of properties
through which an
individual becomes
differentiated. On
one account, the
individual's own
quiddity and its
secondary dispositions
provide a more
integrated explanation
than a simple
reduction to either
a universal essence
or an accidental
bundle. Other
interpretations stress
the different
individuative roles
that Avicenna assigns
to matter, accidents,
and spatiotemporal
determinations.
\citep{BilginTiryakiAvicennaIndividualization2025,PournamdarAvicennaIndividuation2025}

% Structural ID: VS.B01.P01.C08.S13.SS12.P004
\phantomsection\label{para:VS.B01.P01.C08.S13.SS12.P004}

These interpretive
differences matter
for Book I. If
individuation concerns
the determination
of an individual
as distinct from
other individuals,
it cannot simply
be identified with
Identity. Individuation
may be a condition
through which an
individual is
specified, but
the ontological
status of Identity
is not thereby
exhausted. Likewise,
the fact that
an individual can
be distinguished
from another does
not by itself
explain what makes
it the individual
that it is.

% Structural ID: VS.B01.P01.C08.S13.SS12.P005
\phantomsection\label{para:VS.B01.P01.C08.S13.SS12.P005}

The use of perception
in the recognition
or examination of
individuals must
also be kept
distinct from the
ontological constitution
of individuality.
A philosopher may
argue that an
individual is
encountered or
distinguished through
sensible features
without thereby
establishing that
the individual
depends for its
existence or
Identity upon
being perceived.
Book I does not
adopt the cognitive
conditions of
individuation as
the ontological
ground of Identity.


% =============================================================================
% 13. QUIDDITY, INDIVIDUATION, AND IDENTITY
% =============================================================================

\subsection{Quiddity, Individuation, and Identity}
\label{subsec:ch08-avicenna-quiddity-individuation-identity}

% Structural ID: VS.B01.P01.C08.S13.SS13.P001
\phantomsection\label{para:VS.B01.P01.C08.S13.SS13.P001}

The relation between
quiddity and Identity
is the central point
of comparison.
Avicenna's account
of quiddity concerns
the reality or
whatness of a thing.
His account of
individuation concerns
the conditions under
which an individual
is this individual
rather than another.
Book I's account
of Identity concerns
the existent as
the existent that
it is. These
formulations are
closely related,
but their conceptual
functions are not
identical.

% Structural ID: VS.B01.P01.C08.S13.SS13.P002
\phantomsection\label{para:VS.B01.P01.C08.S13.SS13.P002}

The difference must
not be exaggerated.
It would be inadequate
to claim that
quiddity is irrelevant
to Identity or
that individuation
has no ontological
significance. An
existent's whatness
and its particular
constitution may
be indispensable
to an account
of why it is
the existent that
it is. The
question is whether
these dimensions
exhaust Identity
or participate
in its constitution
without being
identical to it.

% Structural ID: VS.B01.P01.C08.S13.SS13.P003
\phantomsection\label{para:VS.B01.P01.C08.S13.SS13.P003}

A further distinction
concerns the difference
between numerical
identity and individuality.
Numerical identity
asks whether an
existent at one
comparison point
is one and the
same existent
as at another.
Individuation asks
how an individual
is specified
as this individual.
Identity in Book I
is not reduced
to either question.
The relations among
these problems
require further
analysis, particularly
when transformation,
persistence, and
possible plurality
of identities are
considered.

% Structural ID: VS.B01.P01.C08.S13.SS13.P004
\phantomsection\label{para:VS.B01.P01.C08.S13.SS13.P004}

Avicenna's account
therefore offers
a substantial
challenge rather
than a ready-made
solution. It
shows that an
ontology of Identity
must explain
the role of
whatness, individual
existence, and
individuative
conditions. Book I
must not merely
repeat its
non-reduction
principle; it
must show how
Identity remains
a distinct
ontological problem
while being related
to the dimensions
through which
an existent
is specified.


% =============================================================================
% 14. POINTS OF CONVERGENCE AND DIVERGENCE
% =============================================================================

\subsection{Points of Convergence and Divergence}
\label{subsec:ch08-avicenna-convergence-divergence}

% Structural ID: VS.B01.P01.C08.S13.SS14.P001
\phantomsection\label{para:VS.B01.P01.C08.S13.SS14.P001}

The first point
of convergence is
the rejection of
an entirely
unexamined notion
of existence.
Avicenna distinguishes
the existent,
the thing,
and the quiddity
through different
conceptual functions.
Book I likewise
requires distinctions
among Being,
Identity,
Determination,
Unity,
and other
related dimensions.
Both approaches
therefore resist
the assumption
that a single
term can perform
every explanatory
task.

% Structural ID: VS.B01.P01.C08.S13.SS14.P002
\phantomsection\label{para:VS.B01.P01.C08.S13.SS14.P002}

A second convergence
concerns the
importance of
causal and
ontological
dependence.
Avicenna offers
a systematic
account of the
dependence of
possible existents
upon causes.
Book I does
not yet adopt
that causal
architecture as
a foundational
principle, but
it must take
seriously the
possibility that
causal relations
may be relevant
to the constitution,
determination,
or manifestation
of existents.

% Structural ID: VS.B01.P01.C08.S13.SS14.P003
\phantomsection\label{para:VS.B01.P01.C08.S13.SS14.P003}

A first divergence
concerns the
central distinction.
Avicenna's principal
metaphysical analysis
concerns essence
and existence,
while Book I's
foundational
distinction concerns
Being and Identity.
A second divergence
concerns the
role of the
First Principle.
Avicenna's Necessary
Existent is
the uncaused
ground of
dependent existents;
Book I has
not identified
its account
of Identity
with a first
causal principle.
A third divergence
concerns the
scope of the
modal account.
Avicenna's distinction
between necessity
and possibility
is central
to his
explanation of
existence, whereas
Book I has
not adopted
a universal
modal reduction
of ontological
relations.

% Structural ID: VS.B01.P01.C08.S13.SS14.P004
\phantomsection\label{para:VS.B01.P01.C08.S13.SS14.P004}

A further divergence
concerns the
status of
individuation.
Avicenna's analyses
of matter,
accidents,
quiddity,
and individual
existence are
directly relevant
to the question
of how an
individual is
specified. Book I,
however, preserves
the distinction
between individuation
and Identity.
It remains open
whether some
individuative
conditions participate
in the constitution
of Identity,
but no
single factor
is adopted
as a universal
criterion.

% Structural ID: VS.B01.P01.C08.S13.SS14.P005
\phantomsection\label{para:VS.B01.P01.C08.S13.SS14.P005}

The comparison
therefore yields
neither simple
agreement nor
simple opposition.
The two
frameworks overlap
in the problems
they investigate,
but their
conceptual
architectures
and explanatory
priorities differ.
Their relation
must be established
through the
arguments and
not through
the superficial
similarity of
translated terms.


% =============================================================================
% 15. THE STRONGEST AVICENNIAN CHALLENGE AND THE RESPONSE OF BOOK I
% =============================================================================

\subsection{The Strongest Avicennian Challenge and the Response of Book I}
\label{subsec:ch08-avicenna-strongest-challenge}

% Structural ID: VS.B01.P01.C08.S13.SS15.P001
\phantomsection\label{para:VS.B01.P01.C08.S13.SS15.P001}

The strongest Avicennian
challenge concerns
the explanatory
necessity of
Identity as
a distinct
ontological dimension.
If quiddity
expresses the
reality by
virtue of which
a thing is
what it is,
and if the
individual is
specified through
the conditions
of its particular
existence,
then what
additional
explanatory work
does the
concept of
Identity perform?
Unless Book I
can answer
this question,
Identity may
appear to
repeat the
function already
performed by
quiddity and
individuation.

% Structural ID: VS.B01.P01.C08.S13.SS15.P002
\phantomsection\label{para:VS.B01.P01.C08.S13.SS15.P002}

The challenge
can be made
more precise.
Book I states
that Identity
is the existent
as the existent
that it is,
but this
formulation may
seem to
restate the
question rather
than answer it.
Avicenna could
argue that
the analysis
of quiddity,
its relation
to existence,
and the
conditions of
individualization
already provide
a more articulated
account of
what a thing
is and
how it
is this
individual. If
Identity adds
no further
distinction,
it risks
becoming an
unexplained
label.

% Structural ID: VS.B01.P01.C08.S13.SS15.P003
\phantomsection\label{para:VS.B01.P01.C08.S13.SS15.P003}

The response
of Book I
must not
consist in
denying the
importance of
quiddity or
individuation.
The response
begins by
distinguishing
the question
of what
a thing is,
the question
of how
an individual
is specified,
and the
question of
the ontological
status of
Identity itself.
These questions
may be related
within the
same existent,
but the
fact that
one contributes
to the
answer of
another does
not establish
their conceptual
equivalence.

% Structural ID: VS.B01.P01.C08.S13.SS15.P004
\phantomsection\label{para:VS.B01.P01.C08.S13.SS15.P004}

The response
also requires
a distinction
between a
condition that
helps explain
an existent
and the
ontological
dimension whose
status is
being explained.
A quiddity
may specify
what a
thing is;
individuative
conditions may
specify how
an individual
is this
individual; and
causal relations
may explain
why it
exists. None
of these
functions,
considered
independently,
has yet
been shown
to exhaust
the ontological
question of
Identity. This
is not
a proof
that Identity
is irreducible
in every
possible sense;
it is
a restriction
on what
can be
inferred from
the Avicennian
analyses alone.

% Structural ID: VS.B01.P01.C08.S13.SS15.P005
\phantomsection\label{para:VS.B01.P01.C08.S13.SS15.P005}

The three
foundational commitments
of Book I
must also
remain separate.
Analytical Distinction
allows Being
and Identity
to be
examined through
different questions.
Ontological Inseparability
rejects the
idea that
an existent
could first
be fully
actual as
an identity-free
entity. Absence
of Ontological
Priority rejects
the assignment
of an
ontological
priority of
one over
the other.
These commitments
do not
deny the
distinction between
quiddity and
existence, nor
do they
automatically settle
the relation
between quiddity
and Identity.

% Structural ID: VS.B01.P01.C08.S13.SS15.P006
\phantomsection\label{para:VS.B01.P01.C08.S13.SS15.P006}

The strongest
response that
can presently
be maintained
is therefore
limited but
substantive: the
Avicennian analysis
of quiddity,
existence,
and individuation
clarifies several
conditions under
which an
existent is
what it
is, but
does not
by itself
demonstrate that
these conditions
are identical
with the
ontological status
of Identity.
Book I
must develop
a positive
account of
this difference
in later
investigations.
Until that
account is
adequately developed,
the relation
between quiddity
and Identity
remains a
serious open
problem rather
than a
settled victory
for either
framework.


% =============================================================================
% 16. THE STRONGEST CHALLENGE FROM BOOK I AND THE AVICENNIAN REPLY
% =============================================================================

\subsection{The Strongest Challenge from Book I and the Avicennian Reply}
\label{subsec:ch08-avicenna-book1-challenge}

% Structural ID: VS.B01.P01.C08.S13.SS16.P001
\phantomsection\label{para:VS.B01.P01.C08.S13.SS16.P001}

The strongest
challenge from
Book I
concerns the
scope of
Avicenna's explanation.
His account
of the
Necessary Existent
and the
causal dependence
of possible
existents explains
why a
contingent existent
exists and
why its
existence is
necessary through
another when
its complete
cause is
present. Yet
the question
remains whether
this causal
and modal
explanation also
provides a
complete account
of the
Identity of
the caused
existent.

% Structural ID: VS.B01.P01.C08.S13.SS16.P002
\phantomsection\label{para:VS.B01.P01.C08.S13.SS16.P002}

The challenge
must not
be formulated
as though
Avicenna ignored
whatness or
individuality. He
explicitly examines
quiddity and
the conditions
under which
individuals are
specified. The
more precise
question is
whether the
causal necessity
of existence,
together with
the account
of quiddity
and individuation,
fully explains
the ontological
status of
Identity, including
the possibility
of multiple
identities,
the relation
between Identity
and internal
plurality, and
the conditions
of Identity
through transformation.

% Structural ID: VS.B01.P01.C08.S13.SS16.P003
\phantomsection\label{para:VS.B01.P01.C08.S13.SS16.P003}

A further
challenge concerns
the difference
between the
uniqueness of
the Necessary
Existent and
the plurality
of caused
existents. The
fact that
the First
Principle is
unique does
not establish
that every
caused existent
possesses exactly
one Identity.
Nor does
the unity
of the
First Principle
explain, without
further argument,
the unity
or plurality
of the
identities of
its effects. A
causal source
and the
identities of
its effects
belong to
different explanatory
questions.

% Structural ID: VS.B01.P01.C08.S13.SS16.P004
\phantomsection\label{para:VS.B01.P01.C08.S13.SS16.P004}

The strongest
Avicennian reply
would be
that the
challenge imposes
a demand
that his
metaphysics does
not leave
unanswered. Quiddity
accounts for
what a
thing is;
the modal
distinction accounts
for the
way its
existence is
related to
its own
essence; the
cause accounts
for its
actual existence;
and individuation
accounts for
the particularity
of the
individual. These
analyses may
be understood
as complementary
rather than
as competing
explanations. The
fact that
Book I
uses a
different foundational
vocabulary does
not establish
that Avicenna's
system lacks
an account
of the
problems it
sets out
to explain.

% Structural ID: VS.B01.P01.C08.S13.SS16.P005
\phantomsection\label{para:VS.B01.P01.C08.S13.SS16.P005}

This reply
is substantial.
Book I
cannot claim
that Avicenna
has failed
merely because
he does
not formulate
the distinction
between Being
and Identity
in the
terms adopted
here. The
comparison must
ask whether
his distinctions
are sufficient
for the
questions he
intends to
answer and
whether they
also resolve
the questions
that Book I
places in
the foreground.
The absence
of a
concept in
one vocabulary
does not
prove the
absence of
the ontological
problem it
might address.

% Structural ID: VS.B01.P01.C08.S13.SS16.P006
\phantomsection\label{para:VS.B01.P01.C08.S13.SS16.P006}

The reply
does not
settle the
comparison. It
shows that
the strongest
criticism from
Book I
must remain
conditional. If
Avicenna's account
of quiddity
and individuation
can be
shown to
explain the
ontological status
of Identity
without reducing
it to
a mere
definition,
a bundle
of determinations,
or a
criterion of
distinction, then
the criticism
must be
revised. Conversely,
if those
analyses leave
an explanatory
gap concerning
Identity, Book I
must specify
that gap
rather than
rely on
terminological
difference alone.


% =============================================================================
% 17. BILATERAL COSTS AND THE SUBSTITUTION TEST
% =============================================================================

\subsection{Bilateral Costs and the Substitution Test}
\label{subsec:ch08-avicenna-bilateral-costs}

% Structural ID: VS.B01.P01.C08.S13.SS17.P001
\phantomsection\label{para:VS.B01.P01.C08.S13.SS17.P001}

The comparison
has a
substantial cost
for Book I.
It cannot
treat quiddity
as a
secondary or
merely linguistic
notion. Avicenna's
analysis demonstrates
that an
account of
Identity must
address the
relation between
what a
thing is,
the conditions
of its
individuality, and
the actuality
through which
it exists. The
ontological status
of Identity
cannot be
developed without
explaining how
these dimensions
are related.

% Structural ID: VS.B01.P01.C08.S13.SS17.P002
\phantomsection\label{para:VS.B01.P01.C08.S13.SS17.P002}

The corresponding
cost for
Avicenna's account,
when examined
from the
questions of
Book I,
is not
that his
system lacks
a theory
of whatness
or individuality.
It is
that the
relation between
quiddity,
individuation,
and Identity
remains a
distinct comparative
question. His
modal and
causal account
explains the
dependence of
possible existents,
but the
extent to
which this
also explains
the full
ontological status
of Identity
requires further
argument.

% Structural ID: VS.B01.P01.C08.S13.SS17.P003
\phantomsection\label{para:VS.B01.P01.C08.S13.SS17.P003}

The substitution
test clarifies
the limits
of equivalence.
Replacing Identity
in Book I
with quiddity
would change
the conceptual
function of
the argument:
quiddity concerns
what a
thing is,
whereas Identity
is the
ontological dimension
that Book I
distinguishes analytically
from Being.
Replacing Being
with existence
in the
Avicennian distinction
would likewise
change the
argument, because
Avicenna's existence
is analysed
in relation
to quiddity
and modality,
not as
a synonym
for the
Being--Identity
relation adopted
here.

% Structural ID: VS.B01.P01.C08.S13.SS17.P004
\phantomsection\label{para:VS.B01.P01.C08.S13.SS17.P004}

Replacing individuation
with Identity
would also
change the
question. Individuation
concerns the
conditions under
which an
individual is
specified as
this individual;
Identity concerns
the existent
as the
existent that
it is. A
relation between
them is
possible and
may be
necessary in
particular accounts,
but their
equivalence has
not been
established. The
same restriction
applies to
Unity, numerical
identity, and
the preservation
of Identity.

% Structural ID: VS.B01.P01.C08.S13.SS17.P005
\phantomsection\label{para:VS.B01.P01.C08.S13.SS17.P005}

The relation
between the
two projects
is therefore
one of
critical proximity
without conceptual
equivalence. Avicenna
offers a
powerful account
of quiddity,
modality, causal
dependence, and
the Necessary
Existent. Book I
offers a
different foundational
analysis of
Being and
Identity. Their
overlap is
philosophically productive,
but neither
system can
be substituted
for the
other without
altering the
questions and
commitments involved.


% =============================================================================
% 18. OPEN PROBLEMS AND CLAIM STATUS
% =============================================================================

\subsection{Open Problems and Claim Status}
\label{subsec:ch08-avicenna-open-problems}

% Structural ID: VS.B01.P01.C08.S13.SS18.P001
\phantomsection\label{para:VS.B01.P01.C08.S13.SS18.P001}

The comparison leaves
several questions
open. First,
what is
the exact
ontological relation
between quiddity
and Identity?
Second, can
the conditions
of individuation
be shown
to participate
in the constitution
of Identity
without being
identical with
it? Third,
does the
modal distinction
between possibility
in itself
and necessity
through another
have implications
for the
conditions under
which Identity
persists through
transformation? Fourth,
can the
uniqueness of
the Necessary
Existent be
related to
the plurality
of identities
without imposing
a universal
one-identity thesis?

% Structural ID: VS.B01.P01.C08.S13.SS18.P002
\phantomsection\label{para:VS.B01.P01.C08.S13.SS18.P002}

A further
question concerns
the status
of causal
dependence. Does
causal dependence
merely explain
why an
existent exists,
or may it
also contribute
to the constitution
of what
that existent
is? If
it contributes
to Identity,
under what
conditions does
that contribution
occur? The
answer cannot
be inferred
from the
fact that
a cause
is necessary
for the
existence of
its effect.
The relation
between causal
dependence and
Identity requires
an independent
argument.

% Structural ID: VS.B01.P01.C08.S13.SS18.P003
\phantomsection\label{para:VS.B01.P01.C08.S13.SS18.P003}

The table below
classifies the
principal Avicennian
claims reconstructed
in this section.
The first
axis records
the functional
or logical
form of
each claim.
The second
axis records
its epistemic
or justificatory
status within
the present
comparative inquiry.
These classifications
are not
attributed to
Avicenna as
his own
taxonomy. They
identify how
the claims
function in
the reconstruction
and how
far the
available arguments
support them.


% =============================================================================
% Project : THE_PHILOSOPHY_OF_VIABILITY
% Book    : Book I — Ontology of Viability
% Book ID : VS.B01
% Part    : Part I — Identity and Manifestation
% Part ID : VS.B01.P01
% Chapter : Chapter 08 — Comparative Philosophical Examination
% Chapter ID : VS.B01.P01.C08
%
% File :
%   4_ASSETS/tables/
%   VS.B01.P01.C08.S13.SS18.P003.TL01.avicenna_claims_and_comparative_limits.tex
%
% Asset :
%   TL01 — Avicenna's Principal Claims and Comparative Limits
%
% Structural Identifier :
%   VS.B01.P01.C08.S13.SS18.P003.TL01
%
% Purpose :
%   To classify the principal Avicennian claims reconstructed
%   in Section 13 and distinguish their logical form from
%   their epistemic or justificatory status in the present
%   comparative examination.
%
% Dependencies :
%   longtable, booktabs, array.
%   These are already loaded by the Book I preamble.
%
% =============================================================================

\begingroup
\small
\renewcommand{\arraystretch}{1.18}
\setlength{\tabcolsep}{4pt}
\setlength{\LTcapwidth}{\linewidth}

% Local table helpers.
\def\colw#1{\dimexpr(\linewidth-8\tabcolsep)*#1/100\relax}

\def\rowsep{\specialrule{0.2pt}{3pt}{3pt}}

\def\tablehead{%
    \toprule
    \textbf{Claim or issue} &
    \textbf{Historical reconstruction} &
    \textbf{Functional or logical form} &
    \textbf{Epistemic or justificatory status} &
    \textbf{Comparative limit for Book I}
    \\
    \midrule
}

\def\grouprow#1{%
    \multicolumn{5}{@{}l}{\footnotesize\textsc{#1}} \\
    \cmidrule{1-5}
    \addlinespace[1pt]
}

\def\EpArg{Argued or Derived Claims}
\def\EpCond{Conditionally Supported Claims}
\def\EpProv{Provisional Constructions}
\def\EpOpen{Unresolved / Open Problems}

\begin{longtable}{@{}
    >{\raggedright\arraybackslash}p{\colw{13}}
    >{\raggedright\arraybackslash}p{\colw{25}}
    >{\raggedright\arraybackslash}p{\colw{14}}
    >{\raggedright\arraybackslash}p{\colw{15}}
    >{\raggedright\arraybackslash}p{\colw{33}}
@{}}

\caption{Avicenna's Principal Claims and Their Comparative Limits}
\label{tab:VS.B01.P01.C08.S13.SS18.P003.TL01}
\\

\tablehead
\endfirsthead

\multicolumn{5}{@{}l}{%
    \footnotesize\textit{
        Table~\thetable\ (continued):
        Avicenna's Principal Claims and Their Comparative Limits
    }
}
\\
\tablehead
\endhead

\multicolumn{5}{@{}r}{%
    \footnotesize\textit{Continued on next page}
}
\\
\endfoot

\bottomrule
\endlastfoot


\grouprow{Metaphysics and the existent}

\textbf{Metaphysics}
&
Metaphysics investigates the existent insofar as it is existent and also establishes the First Principle.
&
Thesis
&
\EpArg
&
The universal subject of metaphysics is not identical with the Necessary Existent or with the totality of causal explanation.
\\
\rowsep

\textbf{The existent}
&
The existent is a primary notion that cannot be defined through a more basic concept.
&
Proposition
&
\EpArg
&
The primacy of the existent does not by itself settle the relation between Being and Identity.
\\
\rowsep

\grouprow{Essence, quiddity, and existence}

\textbf{Essence and existence}
&
The quiddity of a contingent existent does not entail its existence; the relation has modal and causal significance.
&
Thesis
&
\EpArg
&
The distinction is not equivalent to the analytical distinction between Being and Identity.
\\
\rowsep

\textbf{Quiddity in itself}
&
Quiddity considered in itself is neither universal nor particular; these determinations accompany it in modes of existence.
&
Thesis
&
\EpCond
&
The status of quiddity does not establish that it exhausts the ontological constitution of Identity.
\\
\rowsep

\textbf{Thingness and whatness}
&
The reality by virtue of which a thing is what it is is distinguished from the affirmation that the thing exists.
&
Proposition
&
\EpArg
&
The similarity to the formulation of Identity in Book I creates a substantive challenge but not a demonstrated equivalence.
\\
\rowsep

\grouprow{Modality and the First Principle}

\textbf{Necessary Existent}
&
The Necessary Existent is uncaused and necessary in itself.
&
Thesis
&
\EpArg
&
The First Principle is not equivalent to Being as such, Identity, or First Determination.
\\
\rowsep

\textbf{Uniqueness of the Necessary Existent}
&
The Necessary Existent is unique; two independently necessary existents of the same absolute status are not possible within the argument.
&
Proposition
&
\EpArg
&
The uniqueness of the First Principle does not establish that every existent has exactly one identity.
\\
\rowsep

\textbf{Possible existent}
&
A contingent existent is possible in itself and does not contain a sufficient reason for its own existence.
&
Proposition
&
\EpArg
&
Possibility in itself is not a statement of epistemic uncertainty.
\\
\rowsep

\textbf{Necessity through another}
&
A possible existent becomes necessary through its complete cause while remaining possible in itself.
&
Proposition
&
\EpArg
&
Causal necessitation is not equivalent to persistence, invariance, or identity preservation.
\\
\rowsep

\grouprow{Unity, causality, and cosmology}

\textbf{Unity and plurality}
&
Unity is a primary metaphysical notion, but its senses and applications must be distinguished.
&
Proposition
&
\EpArg
&
Unity is not identical with Identity, Singularity, or the unity of an aggregate.
\\
\rowsep

\textbf{Causal series}
&
Caused existents and intermediate causes depend upon a First Principle in the order of causal explanation.
&
Argument
&
\EpCond
&
The argument must not be reduced to a claim about the impossibility of every form of temporal infinity.
\\
\rowsep

\textbf{Emanation}
&
Dependent existents proceed from the First through an ordered causal structure.
&
Thesis
&
\EpCond
&
Emanation is not equivalent to Expression or Manifestation.
\\
\rowsep

\textbf{Cosmological order}
&
Separate intelligences, celestial souls and bodies, and sublunary existents are related through a mediated order.
&
Proposition
&
\EpCond
&
The cosmological order is not a map of the proposed realms and hubs of Book I.
\\
\rowsep

\grouprow{Individuation and the comparative problem}

\textbf{Individuation}
&
Avicenna's discussions involve several possible individuative factors, whose roles differ across contexts and kinds of entities.
&
Open Problem
&
\EpOpen
&
No single principle of individuation is adopted here as a universal criterion of Identity.
\\
\rowsep

\textbf{Quiddity and Identity}
&
The relation between the reality of a thing and the ontological status of Identity requires further analysis.
&
Open Problem
&
\EpOpen
&
Quiddity must not be renamed Identity; its possible constitutive role remains open.
\\
\rowsep

\textbf{Causal dependence and Identity}
&
The causal explanation of the existence of a dependent existent may be relevant to its constitution, but the relation requires independent argument.
&
Hypothesis
&
\EpProv
&
Causal dependence does not automatically provide a complete account of the Identity of the effect.
\\
\rowsep

\textbf{Book I response}
&
Book I distinguishes Being, Identity, quiddity, individuation, and causal dependence without treating them as interchangeable.
&
Thesis
&
\EpProv
&
The positive explanatory role of Identity must be developed beyond a merely negative non-reduction principle.
\\
\rowsep

\textbf{Relation between the systems}
&
The two frameworks address overlapping problems while preserving distinct conceptual and explanatory architectures.
&
Thesis
&
\EpCond
&
The relation is critical proximity without conceptual equivalence or direct doctrinal continuity.
\\

\end{longtable}

\endgroup

% =============================================================================
% END OF TABLE TL01
% =============================================================================


% Structural ID: VS.B01.P01.C08.S13.SS18.P004
\phantomsection\label{para:VS.B01.P01.C08.S13.SS18.P004}

The two
axes must
remain independent.
A proposition
may be
conditionally supported;
a thesis
may be
argued; and
an open
problem may
be precisely
formulated without
being resolved.
The classification
therefore prevents
a historical
reconstruction from
being mistaken
for a
foundational commitment
of Book I.
It also
prevents the
degree of
support for
a claim
from being
confused with
the logical
role that
the claim
performs.

% Structural ID: VS.B01.P01.C08.S13.SS18.P005
\phantomsection\label{para:VS.B01.P01.C08.S13.SS18.P005}

The present
interpretation would
require revision
if further
primary-text analysis
showed that
Avicenna explicitly
identifies quiddity
and Identity
in the
same technical
sense adopted
by Book I,
or if
the different
accounts of
individuation were
shown to
be more
systematically unified
than the
present reconstruction
allows. Conversely,
the comparative
thesis would
be weakened
if Book I
could not
specify the
ontological work
performed by
Identity beyond
the analysis
of quiddity,
determination,
and individuation.

% Structural ID: VS.B01.P01.C08.S13.SS18.P006
\phantomsection\label{para:VS.B01.P01.C08.S13.SS18.P006}

The section
therefore retains
a provisional
comparative conclusion:
Avicenna's metaphysics
offers a
powerful account
of the
relation between
quiddity,
existence,
modality,
and causal
dependence, but
the equivalence
of these
analyses with
the ontological
status of
Identity in
Book I
has not
been established.
The question
remains open
and must
be addressed
in the
subsequent development
of the
ontology rather
than resolved
by terminological
substitution.


% =============================================================================
% 19. TRANSITION TO SUHRAWARDI
% =============================================================================

\subsection{Transition to Suhrawardi}
\label{subsec:ch08-avicenna-transition}

% Structural ID: VS.B01.P01.C08.S13.SS19.P001
\phantomsection\label{para:VS.B01.P01.C08.S13.SS19.P001}

Avicenna's metaphysics
has shown
that the
analysis of
the existent
requires an
account of
quiddity,
existence,
modality,
unity,
causal dependence,
and the
Necessary Existent.
His theory
of individuation
also reveals
that the
question of
what a
thing is
cannot be
separated from
the question
of how
an individual
is specified.
These problems
place a
substantial demand
upon any
ontology of
Identity.

% Structural ID: VS.B01.P01.C08.S13.SS19.P002
\phantomsection\label{para:VS.B01.P01.C08.S13.SS19.P002}

The comparison
has not
established that
quiddity is
Identity, that
existence is
Being in
the sense
adopted by
Book I,
or that
the Necessary
Existent is
First Determination.
It has
instead clarified
the conceptual
and explanatory
differences that
must be
preserved while
examining their
relations. The
remaining problem
is not
to choose
one vocabulary
and discard
the other,
but to
determine what
each account
explains and
what remains
unresolved.

% Structural ID: VS.B01.P01.C08.S13.SS19.P003
\phantomsection\label{para:VS.B01.P01.C08.S13.SS19.P003}

The next
section turns
to Suhrawardi.
It will
examine the
critique of
the primacy
of existence,
the role
of light
in the
Illuminationist
metaphysics,
and the
relation between
quiddity,
presence,
and degrees
of illumination.
The transition
must preserve
the difference
between Avicenna's
modal and
causal account
of existence
and Suhrawardi's
distinctive account
of light and
manifestation. The
latter must
not be
assumed to
resolve the
ontological problem
of Light
as it
will be
investigated in
Part III
of Book I.


% =============================================================================
% END OF SECTION 13 OF CHAPTER 08
% =============================================================================
